% ============================================================
%  Übungsaufgaben Template (my_dhbw Stil, uebung_*.tex)
%  Header "Übungsblatt", grün eingefärbte <details>-Lösungen.
%  Renderer: pdflatex (zweimal)
% ============================================================
\documentclass[11pt, a4paper]{article}

\usepackage[ngerman]{babel}
\usepackage{fontspec}
\setmainfont[
  Path=/usr/share/fonts/TTF/,
  Extension=.ttf,
  BoldFont=DejaVuSans-Bold,
  ItalicFont=DejaVuSans-Oblique,
  BoldItalicFont=DejaVuSans-BoldOblique
]{DejaVuSans}
\setsansfont[
  Path=/usr/share/fonts/TTF/,
  Extension=.ttf,
  BoldFont=DejaVuSans-Bold,
  ItalicFont=DejaVuSans-Oblique,
  BoldItalicFont=DejaVuSans-BoldOblique
]{DejaVuSans}
\setmonofont[Path=/usr/share/fonts/TTF/,Extension=.ttf]{DejaVuSansMono}
\usepackage[top=2cm, bottom=2cm, left=2.4cm, right=2.4cm, footskip=1cm]{geometry}
\usepackage{amsmath, amssymb}
\usepackage{xcolor}
\usepackage{enumitem}
\usepackage{fancyhdr}
\usepackage{lastpage}
\usepackage{tabularx}
\usepackage{booktabs}
\usepackage{microtype}
\usepackage{parskip}

\usepackage{array}
\usepackage{longtable}
\usepackage{ltablex}
\keepXColumns
\usepackage{colortbl}
\usepackage[most,breakable]{tcolorbox}
\usepackage{listings}
\usepackage[normalem]{ulem}
\usepackage{pdflscape}
\usepackage{titlesec}
\usepackage[colorlinks=true, linkcolor=accent, urlcolor=accent]{hyperref}

\renewcommand{\familydefault}{\sfdefault}

% ── Theme ─────────────────────────────────────────────────────
\definecolor{accent}{RGB}{0, 60, 113}
\definecolor{primaryblue}{RGB}{0, 60, 113}
\definecolor{loesung}{RGB}{0, 100, 0}
\definecolor{codebg}{RGB}{245, 245, 245}
\definecolor{figurebg}{RGB}{232, 241, 255}
\definecolor{tablehintbg}{RGB}{240, 255, 240}
\definecolor{solutionbg}{RGB}{240, 252, 240}
\definecolor{rulegray}{RGB}{180, 180, 180}
\definecolor{tableheadbg}{RGB}{0, 60, 113}

% ── Header/Footer ─────────────────────────────────────────────
\pagestyle{fancy}
\fancyhf{}
\renewcommand{\headrulewidth}{0.4pt}
\fancyhead[L]{\footnotesize\color{accent}\nichttitle}
\fancyhead[R]{\footnotesize\color{accent}\nichtsubtitle}
\fancyfoot[R]{\footnotesize\color{accent}Blatt \thepage\,/\,\pageref{LastPage}}

% ── Typografie ────────────────────────────────────────────────
\titleformat{\section}
    {\color{accent}\large\bfseries}{}{0pt}{}
    [\vspace{-4pt}\color{accent}\rule{\linewidth}{1.5pt}]
\titleformat{\subsection}
    {\normalsize\bfseries\color{accent!85!black}}{}{0pt}{}{}
\titleformat{\subsubsection}
    {\small\bfseries\color{accent!70!black}}{}{0pt}{}{}
\titlespacing*{\section}{0pt}{10pt}{3pt}
\titlespacing*{\subsection}{0pt}{6pt}{2pt}
\titlespacing*{\subsubsection}{0pt}{4pt}{1pt}

\setlength{\parindent}{0pt}
\setlength{\parskip}{6pt}
\renewcommand{\arraystretch}{1.2}
\setlist[itemize]{noitemsep, topsep=3pt, parsep=0pt, leftmargin=1.4em}
\setlist[enumerate]{noitemsep, topsep=3pt, parsep=0pt, leftmargin=1.6em}

% ── Boxen: solutionbox = Lösungsbox in grün ──────────────────
\newtcolorbox{figurebox}[1]{
  colback=figurebg, colframe=accent!50,
  title={Abbildung: #1}, fonttitle=\small\bfseries,
  breakable, before skip=6pt, after skip=6pt
}
\newtcolorbox{tablehintbox}[1]{
  colback=tablehintbg, colframe=green!50!black!50,
  title={Tabelle: #1}, fonttitle=\small\bfseries,
  breakable, before skip=6pt, after skip=6pt
}
\newtcolorbox{solutionbox}[1]{
  enhanced,
  colback=solutionbg, colframe=loesung,
  title={#1}, fonttitle=\small\bfseries,
  coltitle=white, colbacktitle=loesung,
  breakable, before skip=8pt, after skip=8pt
}

\lstset{
  basicstyle=\ttfamily\small,
  backgroundcolor=\color{codebg},
  breaklines=true,
  frame=single, framerule=0pt,
  xleftmargin=4pt, xrightmargin=4pt,
  aboveskip=6pt, belowskip=6pt,
  keepspaces=true
}

\providecommand{\nichttitle}{Übungsblatt}
\providecommand{\nichtsubtitle}{}

\renewcommand{\nichttitle}{Analysis 2}
\renewcommand{\nichtsubtitle}{Übungsblatt 1}


\begin{document}

\begin{center}
{\Large\bfseries\color{accent}\nichttitle}
\ifx\nichtsubtitle\empty\else
  \\[2pt]{\normalsize\color{accent!80!black}\nichtsubtitle}
\fi
\\[2pt]
\color{accent}\rule{0.4\linewidth}{0.8pt}\color{black}
\end{center}
\vspace{4pt}

% ── BODY ──────────────────────────────────────────────────────

\section{Übungsblatt 1 — Grundlagen Differenzialrechnung}


\noindent\textcolor{rulegray}{\rule{\linewidth}{0.4pt}}


\paragraph{Aufgabe 1 — Stetigkeit einer stückweisen Funktion \textit{(15 Punkte)}}\mbox{}\\[2pt]

Gegeben sei die Funktion


\[
f(x) = \begin{cases} x^2 + a & \text{für } x < 1 \\ 3x - 2 & \text{für } x \geq 1 \end{cases}
\]


a) Definiere präzise, wann eine Funktion an einer Stelle $x_0$ stetig ist. \textit{(3 P)}

b) Bestimme den Parameter $a \in \mathbb{R}$ so, dass $f$ an der Stelle $x_0 = 1$ stetig ist. \textit{(6 P)}

c) Begründe, ob die Funktion $g(x) = 1/(x-1)$ an der Stelle $x_0 = 1$ stetig sein kann, wenn man $g(1)$ beliebig festlegt. \textit{(6 P)}


\textbf{Lösung:}


a) $f$ ist stetig bei $x_0$, wenn $\lim_{x \to x_0} f(x) = f(x_0)$ gilt. Konkret: linksseitiger Grenzwert = rechtsseitiger Grenzwert = Funktionswert an $x_0$. Anschaulich: kein Sprung, keine Lücke.


b) Linksseitiger Grenzwert bei $x_0 = 1$:

\[
\lim_{x \to 1^-} (x^2 + a) = 1 + a
\]

Rechtsseitiger Grenzwert (= Funktionswert, da $x \geq 1$ den Wert 1 enthält):

\[
\lim_{x \to 1^+} (3x - 2) = 3 \cdot 1 - 2 = 1
\]

Stetigkeitsbedingung: $1 + a = 1 \Rightarrow \boxed{a = 0}$.


c) Für $g(x) = 1/(x-1)$ existiert weder linksseitiger noch rechtsseitiger endlicher Grenzwert bei $x_0 = 1$:

\[
\lim_{x \to 1^-} \frac{1}{x-1} = -\infty, \qquad \lim_{x \to 1^+} \frac{1}{x-1} = +\infty
\]

Da der Grenzwert $\lim_{x \to 1} g(x)$ nicht existiert (analog zum Beispiel $1/x$ bei 0), kann durch keine Festlegung von $g(1)$ Stetigkeit erreicht werden — der zentrale Bestandteil der Definition (Existenz eines beidseitigen, übereinstimmenden Grenzwerts) ist verletzt.



\noindent\textcolor{rulegray}{\rule{\linewidth}{0.4pt}}


\paragraph{Aufgabe 2 — Ableitungsregeln im Schnelldurchgang \textit{(20 Punkte)}}\mbox{}\\[2pt]

Bilde jeweils die erste Ableitung. Gib bei jedem Schritt an, welche Regel(n) du verwendest.


a) $f_1(x) = 5x^4 - 7\ln(x) + 2e^x$ \textit{(4 P)}

b) $f_2(x) = (3x^2 + 1)\cdot \ln(x)$ \textit{(4 P)}

c) $f_3(x) = \dfrac{e^x}{x^2 + 1}$ \textit{(5 P)}

d) $f_4(x) = \ln(4x^3 + 2x)$ \textit{(3 P)}

e) $f_5(x) = e^{x^2 + 3x}$ \textit{(4 P)}


\textbf{Lösung:}


a) Summen- und Potenz-/Logarithmus-/Exponentialregel:

\[
f_1'(x) = 5\cdot 4 x^3 - \frac{7}{x} + 2e^x = 20x^3 - \frac{7}{x} + 2e^x
\]


b) Produktregel mit $g = 3x^2+1$, $h = \ln x$, also $g' = 6x$, $h' = 1/x$:

\[
f_2'(x) = 6x \cdot \ln(x) + (3x^2+1)\cdot \frac{1}{x} = 6x\ln(x) + 3x + \frac{1}{x}
\]


c) Quotientenregel mit $g = e^x$, $h = x^2+1$, $g' = e^x$, $h' = 2x$:

\[
f_3'(x) = \frac{e^x(x^2+1) - e^x\cdot 2x}{(x^2+1)^2} = \frac{e^x(x^2 - 2x + 1)}{(x^2+1)^2} = \frac{e^x(x-1)^2}{(x^2+1)^2}
\]


d) Kettenregel als $(\ln u)' = u'/u$ mit $u = 4x^3 + 2x$, $u' = 12x^2 + 2$:

\[
f_4'(x) = \frac{12x^2 + 2}{4x^3 + 2x}
\]


e) Kettenregel: äußere Funktion $e^{(\cdot)}$, innere $u = x^2 + 3x$, $u' = 2x + 3$:

\[
f_5'(x) = e^{x^2 + 3x}\cdot (2x+3)
\]



\noindent\textcolor{rulegray}{\rule{\linewidth}{0.4pt}}


\paragraph{Aufgabe 3 — Differenzierbarkeit am Knick \textit{(18 Punkte)}}\mbox{}\\[2pt]

Gegeben:


\[
f(x) = \begin{cases} x^2 & \text{für } x \leq 2 \\ b\cdot x + c & \text{für } x > 2 \end{cases}
\]


a) Welche zwei Bedingungen müssen erfüllt sein, damit $f$ an der Stelle $x_0 = 2$ differenzierbar ist? \textit{(4 P)}

b) Bestimme $b, c \in \mathbb{R}$ so, dass $f$ bei $x_0 = 2$ differenzierbar ist. \textit{(8 P)}

c) Ist $f$ mit deinen Werten aus b) auf ganz $\mathbb{R}$ stetig differenzierbar? Begründe. \textit{(6 P)}


\textbf{Lösung:}


a) Differenzierbarkeit bei $x_0$ setzt Stetigkeit voraus. Erforderlich sind also:

\begin{enumerate}
  \item \textbf{Stetigkeit} bei $x_0=2$: linksseitiger Grenzwert = rechtsseitiger Grenzwert = $f(2)$.
  \item \textbf{Übereinstimmende Ableitungen} von links und rechts (kein Knick): $\lim_{x\to 2^-} f'(x) = \lim_{x\to 2^+} f'(x)$.
\end{enumerate}

b) Stetigkeit bei $x_0 = 2$: $2^2 = b\cdot 2 + c \Rightarrow 2b + c = 4 \quad (\text{I})$.


Ableitungen: $(x^2)' = 2x$, also linksseitig $f'(2^-) = 2\cdot 2 = 4$. Rechtsseitig $(bx+c)' = b$, also $f'(2^+) = b$.


Differenzierbarkeit verlangt $b = 4 \quad (\text{II})$.


Einsetzen in (I): $2\cdot 4 + c = 4 \Rightarrow c = -4$.


Ergebnis: $\boxed{b = 4,\ c = -4}$, d.h. rechtsseitig $f(x) = 4x - 4$.


c) Mit diesen Werten gilt:

\begin{itemize}
  \item $f$ ist stetig (durch (I) konstruiert) und differenzierbar an $x_0=2$ (durch (II)).
  \item Auf $(-\infty, 2)$ ist $f'(x) = 2x$ — als Polynom stetig.
  \item Auf $(2, \infty)$ ist $f'(x) = 4$ — konstant, also stetig.
  \item Bei $x_0 = 2$: $\lim_{x\to 2^-} 2x = 4$ und $\lim_{x\to 2^+} 4 = 4 = f'(2)$, $f'$ also stetig.
\end{itemize}

Damit ist $f$ auf ganz $\mathbb{R}$ \textbf{stetig differenzierbar}.



\noindent\textcolor{rulegray}{\rule{\linewidth}{0.4pt}}


\paragraph{Aufgabe 4 — Fehler finden in der Ableitung \textit{(12 Punkte)}}\mbox{}\\[2pt]

Eine Studierende rechnet:


\begin{quote}
\textbf{Aufgabe:} Bestimme $f'(x)$ für $f(x) = \dfrac{\ln(x)}{x^3}$.
\end{quote}
>

\begin{quote}
\textbf{Lösung:}
\end{quote}
\begin{quote}
Mit Quotientenregel und $g = \ln(x)$, $h = x^3$, $g' = 1/x$, $h' = 3x^2$:
\end{quote}
>

\begin{quote}
$f'(x) = \dfrac{g'\cdot h - g\cdot h'}{h} = \dfrac{\frac{1}{x}\cdot x^3 - \ln(x)\cdot 3x^2}{x^3} = \dfrac{x^2 - 3x^2\ln(x)}{x^3} = \dfrac{1 - 3\ln(x)}{x}$
\end{quote}

a) Identifiziere den Fehler in der Anwendung der Quotientenregel. \textit{(4 P)}

b) Führe die korrekte Rechnung durch und vereinfache so weit wie möglich. \textit{(8 P)}


\textbf{Lösung:}


a) Die Quotientenregel lautet

\[
\left(\frac{g}{h}\right)' = \frac{g'h - gh'}{h^2}
\]

Im Nenner muss $h^2$ stehen, nicht $h$. Die Studierende hat im Nenner $x^3$ statt $(x^3)^2 = x^6$ verwendet — der Faktor $x^3$ im Nenner fehlt.


b) Korrekt:

\[
f'(x) = \frac{\frac{1}{x}\cdot x^3 - \ln(x)\cdot 3x^2}{(x^3)^2} = \frac{x^2 - 3x^2\ln(x)}{x^6} = \frac{x^2(1 - 3\ln(x))}{x^6} = \boxed{\frac{1 - 3\ln(x)}{x^4}}
\]


(Ergebnis der Studierenden $\frac{1-3\ln x}{x}$ unterscheidet sich um Faktor $x^3$ vom korrekten Wert.)



\noindent\textcolor{rulegray}{\rule{\linewidth}{0.4pt}}


\paragraph{Aufgabe 5 — Vergleich: Produkt- vs. Quotientenregel \textit{(15 Punkte)}}\mbox{}\\[2pt]

Gegeben: $f(x) = \dfrac{x^2 + 1}{e^x}$.


a) Berechne $f'(x)$ mit der \textbf{Quotientenregel}. \textit{(5 P)}

b) Schreibe $f$ um als $f(x) = (x^2+1)\cdot e^{-x}$ und berechne $f'(x)$ mit \textbf{Produkt- und Kettenregel}. \textit{(6 P)}

c) Zeige durch Umformung, dass beide Ergebnisse übereinstimmen, und nenne ein Kriterium, wann du in der Praxis welchen Weg bevorzugen würdest. \textit{(4 P)}


\textbf{Lösung:}


a) Quotientenregel mit $g = x^2+1$, $h = e^x$, $g' = 2x$, $h' = e^x$:

\[
f'(x) = \frac{2x\cdot e^x - (x^2+1)\cdot e^x}{(e^x)^2} = \frac{e^x(2x - x^2 - 1)}{e^{2x}} = \frac{2x - x^2 - 1}{e^x}
\]


b) Produktregel mit $g = x^2+1$, $h = e^{-x}$. Ableitung von $h$ via Kettenregel: $(e^{-x})' = e^{-x}\cdot(-1) = -e^{-x}$. Außerdem $g' = 2x$.

\[
f'(x) = 2x\cdot e^{-x} + (x^2+1)\cdot(-e^{-x}) = e^{-x}(2x - x^2 - 1)
\]


c) Wegen $e^{-x} = 1/e^x$ ist

\[
e^{-x}(2x - x^2 - 1) = \frac{2x - x^2 - 1}{e^x}
\]

also identisch zu (a).


\textbf{Praxis-Kriterium:} Wenn der Nenner eine "einfache" Funktion ist, deren Reziprokes / negative Hochzahl bequem zu schreiben ist (wie $e^x \to e^{-x}$ oder $x^n \to x^{-n}$), ist Umformen auf Produkt + Kettenregel meist rechenökonomischer und vermeidet das Quadrat im Nenner. Bei "echten" Brüchen wie $\frac{\ln x}{x^2+1}$, wo der Nenner kein einfaches Potenzgesetz erlaubt, ist die Quotientenregel direkt und übersichtlicher.



\noindent\textcolor{rulegray}{\rule{\linewidth}{0.4pt}}


\paragraph{Aufgabe 6 — Transfer: Tangente an einen Kostenverlauf \textit{(20 Punkte)}}\mbox{}\\[2pt]

Ein Server-Lastmodell beschreibt die durchschnittliche Antwortzeit (in ms) in Abhängigkeit von der Anzahl gleichzeitiger Anfragen $x > 0$ durch


\[
T(x) = 2\cdot \ln(x) + \frac{1}{50}x^2
\]


a) Berechne $T'(x)$ und interpretiere die Bedeutung von $T'(x_0)$ im Kontext des Modells. \textit{(5 P)}

b) Bestimme die Gleichung der Tangente an $T$ an der Stelle $x_0 = 10$ in der Form $y = m\cdot x + n$. \textit{(8 P)}

c) Begründe rein mit den Begriffen aus dem Kapitel, warum $T$ auf ganz $(0, \infty)$ differenzierbar ist, an $x = 0$ aber \textbf{nicht einmal stetig fortsetzbar} ist. \textit{(7 P)}


\textbf{Lösung:}


a) Mit Logarithmus- und Potenzregel sowie Summenregel:

\[
T'(x) = \frac{2}{x} + \frac{2x}{50} = \frac{2}{x} + \frac{x}{25}
\]


Interpretation: $T'(x_0)$ ist die momentane Änderungsrate der Antwortzeit. Bei $x_0$ gleichzeitigen Anfragen erhöht sich die Antwortzeit pro zusätzlicher Anfrage näherungsweise um $T'(x_0)$ ms (lokale Steigung des Antwortzeit-Verlaufs).


b) Tangentensteigung:

\[
m = T'(10) = \frac{2}{10} + \frac{10}{25} = 0{,}2 + 0{,}4 = 0{,}6
\]


Funktionswert:

\[
T(10) = 2\ln(10) + \frac{100}{50} = 2\ln(10) + 2
\]


Tangentengleichung über Punkt-Steigungs-Form $y - T(10) = m(x - 10)$:

\[
y = 0{,}6(x - 10) + 2\ln(10) + 2 = 0{,}6\, x - 6 + 2\ln(10) + 2 = 0{,}6\,x + (2\ln(10) - 4)
\]


Numerisch: $2\ln(10) \approx 4{,}6052$, also $n \approx 0{,}6052$.


\[
\boxed{y = 0{,}6\,x + (2\ln(10) - 4) \approx 0{,}6\,x + 0{,}6052}
\]


c) Auf $(0, \infty)$:

\begin{itemize}
  \item $\ln(x)$ ist dort differenzierbar mit $(\ln x)' = 1/x$.
  \item $\frac{1}{50}x^2$ ist als Polynom überall differenzierbar.
  \item Die Summe differenzierbarer Funktionen ist (Summenregel) wieder differenzierbar.
\end{itemize}

Folglich ist $T$ auf $(0, \infty)$ differenzierbar und damit dort insbesondere stetig.


An $x_0 = 0$: Differenzierbarkeit setzt Stetigkeit voraus. Es gilt jedoch

\[
\lim_{x\to 0^+} \ln(x) = -\infty
\]

und damit

\[
\lim_{x\to 0^+} T(x) = -\infty + 0 = -\infty
\]

Es existiert also kein endlicher rechtsseitiger Grenzwert, der gleich einem Funktionswert $T(0)$ sein könnte — analog zum Kapitel-Beispiel "$1/x$ ist nicht stetig bei 0". Daher lässt sich $T$ bei 0 nicht stetig fortsetzen, und Differenzierbarkeit an dieser Stelle ist erst recht unmöglich.





% ──────────────────────────────────────────────────────────────

\end{document}
